\subsection{The connecting homomorphism and the exact homology sequence}

In this section, we consider an exact sequence of complexes

\[
0 \longrightarrow \mathrm{C} ^ {\prime} \stackrel {u} {\longrightarrow} \mathrm{C} \stackrel {v} {\longrightarrow} \mathrm{C} ^ {\prime \prime} \longrightarrow 0;
\]

denote by the same letter d the differentials of C, C' and C''\,'.

Let \(\Gamma\) be the set of \(x \in \mathbf{C}\) such that \(dx \in \operatorname{Im}(u)\) ; for \(x \in \Gamma\) , we have

\[
d \left(^ {- 1} u ^ {1} (d x)\right) = ^ {- 1} u ^ {1} (d d (x)) = 0,
\]

hence \(\overline{u}^{1}(dx)\in\mathbf{Z}(\mathbf{C}^{\prime})\) ; we also have \(dv(x)=v(dx)\in\operatorname{Im}(v\circ u)=0\) , hence \(v(x)\in\mathbf{Z}(\mathbf{C}^{\prime\prime})\) ; consider then the linear map \(\varphi:\Gamma\to\mathrm{H}(\mathbf{C}^{\prime\prime})\times\mathrm{H}(\mathbf{C}^{\prime})\) which maps every element \(x\in\Gamma\) to the class of \((v(x),\overline{u}^{1}(dx))\) .

Lemma 1. --- The image \(\varphi(\Gamma)\) of \(\Gamma\) in \(\mathrm{H}(\mathbf{C}^{\prime\prime}) \times \mathrm{H}(\mathbf{C}^{\prime})\) is the graph of a graded A-homomorphism of degree -1 of \(\mathrm{H}(\mathbf{C}^{\prime\prime})\) to \(\mathrm{H}(\mathbf{C}^{\prime})\) .

\begin{enumerate}
\def\labelenumi{\alph{enumi})}
\item
  If \(x \in \Gamma\) and if \(v(x) \in \mathbf{B}(\mathbf{C}'')\) , then \(\bar{u}^1(dx) \in \mathbf{B}(\mathbf{C}')\) : there exists indeed \(z'' \in \mathbf{C}''\) such that \(v(x) = dz''\) , then \(z \in \mathbf{C}\) such that \(z'' = v(z)\) , hence \(v(x) = v(dz)\) , then \(t' \in \mathbf{C}'\) such that \(x - dz = u(t')\) , which gives \(dx = u(dt')\) , hence \(\bar{u}^1(dx) = dt' \in \mathbf{B}(\mathbf{C}')\) .
\item
  Every element of \(\mathbf{Z}(\mathbf{C}^{\prime \prime})\) is the image under \(v\) of an element \(x\) of \(\mathbf{C}\) such that \(v(dx) = 0\) , i.e.~\(dx \in \operatorname{Im} u\) , that is, such that \(x \in \Gamma\) .
\item
  It follows from a) and b) that \(\varphi(\Gamma)\) is indeed a functional graph; since \(\varphi\) is bihomogeneous of bidegree \((0, -1)\) , this completes the proof.
\end{enumerate}

The graded homomorphism of degree -1 of H(C'') in H(C') thus defined is called the connecting homomorphism relative to the exact sequence (u, v); it is denoted \(\partial(u, v)\) or \(\partial_{u,v}\) or simply \(\partial\) . Its homogeneous components are denoted

\[
\partial_ {n} (u, v): \mathrm{H} _ {n} (\mathbf {C} ^ {\prime \prime}) \rightarrow \mathrm{H} _ {n - 1} (\mathbf {C} ^ {\prime}) \quad \text { and } \quad \partial^ {n} (u, v): \mathrm{H} ^ {n} (\mathbf {C} ^ {\prime \prime}) \rightarrow \mathrm{H} ^ {n + 1} (\mathbf {C} ^ {\prime}).
\]

By definition, to construct the image of a class \(\alpha\in\mathrm{H}_{n}(\mathbf{C}^{\prime\prime})\) by \(\partial\) we choose a cycle \(z^{\prime\prime}\in\mathbf{Z}_{n}(\mathbf{C}^{\prime\prime})\) in the class \(\alpha\) , then an element \(x\) of \(\mathbf{C}_{n}\) such that \(v(x)=z^{\prime\prime}\) ; then \(dx\) is of the form \(u(t')\) , \(t^{\prime}\in\mathbf{C}_{n-1}^{\prime}\) , and \(\partial(\alpha)\) is the homology class of \(t'\) .

In terms of correspondences, \(\partial_{n}(u,v)\) is therefore obtained from the correspondence \(u_{n-1}^{-1}\circ d_{n}\circ v_{n}^{-1}\) between \(\mathbf{C}_{n}^{\prime\prime}\) and \(\mathbf{C}_{n-1}^{\prime}\) , by passing to subsets \(\mathbf{Z}_{n}(\mathbf{C}^{\prime\prime})\) and \(\mathbf{Z}_{n-1}(\mathbf{C}^{\prime})\) , then to their quotients \(\mathbf{H}_{n}(\mathbf{C}^{\prime\prime})\) and \(\mathbf{H}_{n-1}(\mathbf{C}^{\prime})\) . This shows in particular that, if one replaces \(\mathbf{C},\mathbf{C}^{\prime},\mathbf{C}^{\prime\prime},u,v\) by \(\mathbf{C}(p),\mathbf{C}^{\prime}(p),\mathbf{C}^{\prime\prime}(p),u(p),v(p)\) , we have

\[
\partial (u (p), v (p)) = (- 1) ^ {p} \partial (u, v);\tag{4}
\]

likewise, if \(\lambda\) and \(\mu\) are two invertible elements of the center of A, one has

\[
\partial (\lambda u, \mu v) = \lambda^ {- 1} \mu^ {- 1} \partial (u, v).\tag{5}
\]

One can also relate \(\partial(u,v)\) to the snake diagram (X, p.~4). According to loc. cit., prop. 2, the sequences

\[
0 \longrightarrow Z _ {n} (C ^ {\prime}) \xrightarrow {Z _ {n} (u)} Z _ {n} (C) \xrightarrow {Z _ {n} (v)} Z _ {n} (C ^ {\prime \prime})
\]

and

\[
\mathrm{C} _ {n} ^ {\prime} / \mathrm{B} _ {n} (\mathrm{C} ^ {\prime}) \xrightarrow {\overline {{{u}}} _ {n}} \mathrm{C} _ {n} / \mathrm{B} _ {n} (\mathrm{C}) \xrightarrow {\overline {{{v}}} _ {n}} \mathrm{C} _ {n} ^ {\prime \prime} / \mathrm{B} _ {n} (\mathrm{C} ^ {\prime}) \longrightarrow 0,
\]

where \(\overline{u}_{n}\) and \(\overline{v}_{n}\) are deduced from \(u\) and \(v\) , are exact. Using the canonical exact sequences \((\mathrm{V}_{n})\) , one obtains a commutative diagram with exact rows and columns

\[
\begin{array}{c c c c} 0 & 0 & 0 \\ \downarrow & \mathrm{H} _ {n} (\mathrm{C} ^ {\prime}) & \xrightarrow {\mathrm{H} _ {n} (u)} & \mathrm{H} _ {n} (\mathrm{C}) \\ \downarrow & \mathrm{C} _ {n} ^ {\prime} / \mathrm{B} _ {n} (\mathrm{C} ^ {\prime}) & \xrightarrow {\bar {u} _ {n}} & \mathrm{C} _ {n} / \mathrm{B} _ {n} (\mathrm{C}) \\ \downarrow & 0 \longrightarrow Z _ {n - 1} (\mathrm{C} ^ {\prime}) & \xrightarrow {Z _ {n - 1} (u)} & Z _ {n - 1} (\mathrm{C}) \\ \downarrow & H _ {n - 1} (\mathrm{C} ^ {\prime}) & \xrightarrow {\mathrm{H} _ {n - 1} (u)} & H _ {n - 1} (\mathrm{C}) \\ \downarrow & 0 & 0 & 0 \end{array}
\]

The homomorphism \(\mathrm{H}_{n}(\mathbf{C}^{\prime \prime})\to \mathrm{H}_{n - 1}(\mathbf{C}^{\prime})\) associated with this diagram (loc. cit., prop. 2, (iii)) coincides by construction with \(\partial_n(u,v)\) . This also implies that the sequence of homomorphisms \((\mathrm{H}_n(u),\mathrm{H}_n(v),\partial_n(u,v),\mathrm{H}_{n - 1}(u),\mathrm{H}_{n - 1}(v))\) is exact; consequently:

THEOREM 1. --- The infinite sequence of homomorphisms of A-modules

\[
\begin{array}{r l} \dots & \to \mathrm{H} _ {n + 1} (\mathrm{C} ^ {\prime \prime}) \xrightarrow {\partial_ {n + 1} (u , v)} \mathrm{H} _ {n} (\mathrm{C} ^ {\prime}) \xrightarrow {\mathrm{H} _ {n} (u)} \mathrm{H} _ {n} (\mathrm{C}) \xrightarrow {\mathrm{H} _ {n} (v)} \mathrm{H} _ {n} (\mathrm{C} ^ {\prime \prime}) \\ & \xrightarrow {\partial_ {n} (u , v)} \mathrm{H} _ {n - 1} (\mathrm{C} ^ {\prime}) \xrightarrow {\mathrm{H} _ {n - 1} (u)} \mathrm{H} _ {n - 1} (\mathrm{C}) \xrightarrow {\mathrm{H} _ {n - 1} (v)} \mathrm{H} _ {n - 1} (\mathrm{C} ^ {\prime \prime}) \xrightarrow {\partial_ {n - 1} (u , v)} \mathrm{H} _ {n - 2} (\mathrm{C} ^ {\prime}) \to \dots \end{array}
\]

is exact.

This sequence is called the exact homology sequence associated with the exact sequence \((u, v)\) ; it is sometimes written in the form of an exact triangle of A-modules

\[
\begin{array}{c} \mathrm{H(C)} \\ \mathrm{H(u)} \nearrow \searrow \mathrm{H(v)} \\ \mathrm {H(C^ {\prime})} \xleftarrow [ \partial (u , v) ]{} \mathrm {H(C^{\prime\prime})}. \end{array}
\]

COROLLARY 1. --- If two of the complexes C, C', C'' have zero homology, then the third also has zero homology. For u (resp. v) to be a homologism, it is necessary and sufficient that C'' (resp. C') have zero homology. For \(\partial(u, v)\) to be bijective, it is necessary and sufficient that C have zero homology.

COROLLARY 2. --- Let u be a morphism of complexes. If Ker u and Coker u have zero homology, then u is a homologism.

Indeed, let \(u: \mathrm{E} \to \mathrm{E}'\) a morphism of complexes. If Ker \(u\) (resp. Coker \(u\) ) has zero homology, then the canonical morphism \(\mathrm{E} \to \operatorname{Im} u\) (resp. \(\operatorname{Im} u \to \mathrm{E}'\) ) is a homologism by Corollary 1.

PROPOSITION 2. — Consider a commutative diagram of complexes with exact rows

\[
\begin{array}{c} 0 \longrightarrow \mathrm{C} ^ {\prime} \stackrel {{u}} {{\longrightarrow}} \mathrm{C} \stackrel {{v}} {{\longrightarrow}} \mathrm{C} ^ {\prime \prime} \longrightarrow 0 \\ f ^ {\prime} \Big \downarrow \qquad \qquad \qquad f \Big \downarrow \qquad \qquad f ^ {\prime \prime} \Big \downarrow \\ 0 \longrightarrow \mathrm{C} _ {1} ^ {\prime} \stackrel {{u _ {1}}} {{\longrightarrow}} \mathrm{C} _ {1} \stackrel {{v _ {1}}} {{\longrightarrow}} \mathrm{C} _ {1} ^ {\prime \prime} \longrightarrow 0. \end{array}
\]

Then \(\mathbf{H}(f')\circ \partial (u,v) = \partial (u_1,v_1)\circ \mathbf{H}(f'')\)

Let \(\alpha'' \in \mathrm{H}(\mathbf{C}'')\) ; let \(z''\) be a cycle of class \(\alpha''\) and \(x\) an element of C such that \(v(x) = z''\) ; we have

\[
\begin{array}{r l} \big (\partial_ {u _ {1}, v _ {1}} \circ \mathrm{H} (f ^ {\prime \prime}) \big) (\alpha^ {\prime \prime}) = & \partial_ {u _ {1}, v _ {1}} (\overline {{f ^ {\prime \prime} (z ^ {\prime \prime})}}) = \overline {{u _ {1} ^ {- 1} (d f (x))}} = \overline {{f ^ {\prime} (u ^ {- 1} (d x))}} = \\ & = \mathrm{H} (f ^ {\prime}) (\overline {{u ^ {- 1} (d x)}}) = (\mathrm{H} (f ^ {\prime}) \circ \partial_ {u, v}) (\alpha^ {\prime \prime}). \end{array}
\]

Example. — Let C be a complex. Consider the canonical exact sequence

\[
0 \longrightarrow \mathrm{Z} (\mathrm{C}) \stackrel {j} {\longrightarrow} \mathrm{C} \stackrel {\delta} {\longrightarrow} \mathrm{B} (\mathrm{C}) (- 1) \longrightarrow 0,\tag{I}
\]

and let \(i: \mathbf{B}(\mathbf{C}) \to \mathbf{Z}(\mathbf{C})\) be the canonical injection. Then the connecting homomorphism \(\partial(j, \delta): \mathrm{H}(\mathrm{B}(\mathrm{C})) (-1) \to \mathrm{H}(\mathrm{Z}(\mathrm{C})) (-1)\) is identified with \(\mathrm{H}(i) (-1)\) , as one verifies immediately. Since \(\mathrm{H}(\delta) = 0\) , the exact homology sequence associated with (I) decomposes into the short exact sequences

\[
0 \longrightarrow \mathrm{B} _ {n} (\mathrm{C}) \stackrel {i} {\longrightarrow} Z _ {n} (\mathrm{C}) \stackrel {\mathrm{H} (j)} {\longrightarrow} \mathrm{H} _ {n} (\mathrm{C}) \longrightarrow 0.\tag{\((\mathrm{II}_{\mathfrak{n}})\)}
\]

* Applications:

\begin{enumerate}
\def\labelenumi{\arabic{enumi})}
\tightlist
\item
  Singular homology
\end{enumerate}

Let A be a ring. For every topological space X, one defines the singular complex C(X, A) of X with coefficients in A as follows:

In \(\mathbf{R}^{(N)}\) , let \((e_n)\) be the canonical basis; one calls canonical \(n\) -simplex the convex hull \(\Delta_n\) of \(\{e_0, \ldots, e_n\}\) . For \(i \in \{0, \ldots, n\}\) , one defines the affine map \(\iota_i: \Delta_{n-1} \to \Delta_n\) by \(\iota_i(e_k) = e_k\) for \(k < i\) and \(\iota_i(e_k) = e_{k+1}\) for \(k \geqslant i\) . One denotes by \(\mathbf{C}_n(\mathbf{X}, \mathbf{A})\) the A-module \(\mathbf{A}^{(\Sigma_n(\mathbf{X}))}\) , where \(\Sigma_n(\mathbf{X})\) is the set of continuous maps from \(\Delta_n\) to \(\mathbf{X}\) ; for \(n < 0\) , we set \(\mathbf{C}_n = 0\) . For \(i \in \{0, \ldots, n\}\) , one defines the linear map \(\partial_{n,i}: \mathbf{C}_n(\mathbf{X}, \mathbf{A}) \to \mathbf{C}_{n-1}(\mathbf{X}, \mathbf{A})\) by \(\partial_{n,i}(e_s) = e_{s \circ \iota_i}\) for \(s \in \Sigma_n(\mathbf{X})\) , and one sets \(d_n = \Sigma (-1)^t \partial_{n,i}\) . One verifies that

\[
\dots \mathrm{C} _ {n} (\mathrm{X}, \mathrm{A}) \xrightarrow {d _ {n}} \mathrm{C} _ {n - 1} (\mathrm{X}, \mathrm{A}) \rightarrow \dots
\]

is a complex. Its homology is called the singular homology of X with coefficients in A and is denoted H(X, A) or simply H(X).

If Y is a subspace of X, one denotes by C(X, Y, A) the quotient complex C(X, A)/C(Y, A), and by H(X, Y, A) its homology. It follows from Theorem I that one has an exact sequence:

\[
\dots \rightarrow \mathrm{H} _ {n} (\mathrm{Y}, \mathrm{A}) \rightarrow \mathrm{H} _ {n} (\mathrm{X}, \mathrm{A}) \rightarrow \mathrm{H} _ {n} (\mathrm{X}, \mathrm{Y}, \mathrm{A}) \rightarrow \mathrm{H} _ {n - 1} (\mathrm{Y}, \mathrm{A}) \rightarrow \mathrm{H} _ {n - 1} (\mathrm{X}, \mathrm{A}) \rightarrow \dots
\]

\begin{enumerate}
\def\labelenumi{\arabic{enumi})}
\setcounter{enumi}{1}
\tightlist
\item
  Finite cellular complexes
\end{enumerate}

Let X be a Hausdorff topological space. A finite cellular decomposition of X is given by an increasing sequence \((X_{n})_{n\in\mathbb{Z}}\) of closed subspaces of X satisfying the following conditions: (i) \(X_{n} = \varnothing\) for n \textless{} 0 ;

\begin{enumerate}
\def\labelenumi{(\roman{enumi})}
\setcounter{enumi}{1}
\item
  there exists an N such that \(X_{N} = X\) (hence \(X_{n} = X\) for n \textgreater{} N);
\item
  for all \(n\) , the space \(\mathbf{X}_{n} - \mathbf{X}_{n-1}\) has only finitely many connected components, called cells of dimension \(n\) ;
\item
  for every n, and for every connected component C of \(X_{n} - X_{n-1}\) there exists a homeomorphism of the open Euclidean ball \(\dot{B}_{n}\) of dimension n (TG, VI, p.~10) onto C which extends to a continuous map from the closed ball into X.
\end{enumerate}

One can show that these conditions imply that \(\mathbf{H}_n(\mathbf{X}_n, \mathbf{X}_{n-1}, \mathbf{A})\) is a free A-module \(\Gamma_n\) of rank equal to the number of cells of dimension \(n\) , and that \(\mathbf{H}_i(\mathbf{X}_n, \mathbf{X}_{n-1}, \mathbf{A}) = 0\) for \(i \neq n\) . We have \(\mathbf{C}(\mathbf{X}_n, \mathbf{X}_{n-1}, \mathbf{A}) = \mathbf{C}(\mathbf{X}_n, \mathbf{X}_{n-2}, \mathbf{A}) / \mathbf{C}(\mathbf{X}_{n-1}, \mathbf{X}_{n-2}, \mathbf{A})\) whence an exact sequence

\[
\mathrm{H} _ {n} \left(\mathrm{X} _ {n}, \mathrm{X} _ {n - 2}\right) \longrightarrow \mathrm{H} _ {n} \left(\mathrm{X} _ {n}, \mathrm{X} _ {n - 1}\right) \xrightarrow {d _ {n}} \mathrm{H} _ {n - 1} \left(\mathrm{X} _ {n - 1}, \mathrm{X} _ {n - 2}\right) \longrightarrow \mathrm{H} _ {n - 1} \left(\mathrm{X} _ {n}, \mathrm{X} _ {n - 2}\right),
\]

with a connecting homomorphism \(d_{n}:\Gamma_{n}\to\Gamma_{n-1}\) . We have \(d_{n}\circ d_{n+1}=0\) , which allows one to define a complex \(\Gamma:\cdots\to\Gamma_{n}\xrightarrow{d_{n}}\Gamma_{n-1}\to\cdots\) .

The exact sequence

\[
\mathrm{H} _ {n + 1} \left(\mathrm{X} _ {p}, \mathrm{X} _ {p - 1}\right)\rightarrow \mathrm{H} _ {n} \left(\mathrm{X} _ {p - 1}\right)\rightarrow \mathrm{H} _ {n} \left(\mathrm{X} _ {p}\right)\rightarrow \mathrm{H} _ {n} \left(\mathrm{X} _ {p}, \mathrm{X} _ {p - 1}\right)\rightarrow \mathrm{H} _ {n - 1} \left(\mathrm{X} _ {p - 1}\right)
\]

shows by induction on \(p\) that \(\mathrm{H}_{n}(\mathrm{X}_{p})=0\) for \(p<n\) , that \(\mathrm{H}_{n}(\mathrm{X}_{n})=\mathrm{Ker}\left(d_{n}:\Gamma_{n}\to\Gamma_{n-1}\right)\) and that \(\mathrm{H}_{n}(\mathrm{X}_{p})=\mathrm{H}_{n}(\Gamma)\) for \(p>n\) In particular \(\mathrm{H}_{n}(\mathrm{X})\) is identified with \(\mathrm{H}_{n}(\Gamma)\) .

Example. --- Consider the product of spheres \(\mathbf{S}_{2} \times \mathbf{S}_{2}\) and the complex projective space \(\mathbf{P}_{2}(\mathbf{C})\) . Let \(b \in \mathbf{S}_{2}\) ; one defines a cell decomposition \((\mathrm{Y}_{n})\) of \(\mathrm{Y} = \mathbf{S}_{2} \times \mathbf{S}_{2}\) by setting

\[
\mathrm{Y} _ {0} = \mathrm{Y} _ {1} = \{(b, b) \}, \quad \mathrm{Y} _ {2} = \mathrm{Y} _ {3} = (\{b \} \times \mathbf {S} _ {2}) \cup (\mathbf {S} _ {2} \times \{b \}) \quad \text {and} \quad \mathrm{Y} _ {4} = \mathbf {S} _ {2} \times \mathbf {S} _ {2};
\]

this decomposition has one cell of dimension 0, two of dimension 2, and one of dimension 4. The differentials of the associated complex are necessarily zero, hence \(\mathrm{H}_{0}(\mathrm{Y})\) , \(\mathrm{H}_{2}(\mathrm{Y})\) and \(\mathrm{H}_{4}(\mathrm{Y})\) are free of ranks 1, 2, and 1 respectively, and \(\mathrm{H}_{n}(\mathrm{Y}) = 0\) for \(n \notin \{0, 2, 4\}\) .

We obtain a cellular decomposition \((\mathbf{Z}_{n})\) of \(\mathbf{P}_{2}(\mathbf{C})\) by setting

\[
\mathbf {Z} _ {0} = \mathbf {Z} _ {1} = \{c \}, \quad \mathbf {Z} _ {2} = \mathbf {Z} _ {3} = \mathbf {P} _ {1} (\mathbf {C}), \quad \mathbf {Z} _ {4} = \mathbf {P} _ {2} (\mathbf {C}),
\]

the space \(\mathbf{P}_{1}(\mathbf{C})\) being embedded in \(\mathbf{P}_{2}(\mathbf{C})\) (TG, VIII, p.~20), and \(c\) being a point of \(\mathbf{P}_{1}(\mathbf{C})\) ; this decomposition has one cell of dimension 0, one of dimension 2, and one of dimension 4. Here again the differentials of the complex are necessarily zero, and it follows that \(\mathbf{H}_{n}(\mathbf{P}_{2}(\mathbf{C}))\) is isomorphic to A for \(n \in \{0, 2, 4\}\) and 0 otherwise.

Since the homology modules in degree 2 of the two spaces under consideration are free of rank 2 and 1, respectively, these spaces are not homeomorphic. *

